\documentclass[11pt,a4paper]{article}

\def\courseTitle{Industrial Security (Master)}
\def\labNumber{05}
\def\labTitle{Secure-by-Design Architecture Workshop}
\def\labSection{05}
\def\labTime{4 hours, group of 3--4}

% =====================================================================
% Shared preamble for lab sheets.
%
% Caller must \documentclass{article} first, then define:
%   \def\courseTitle{...}
%   \def\labNumber{NN}
%   \def\labTitle{...}
%   \def\labTime{e.g. "30--45 min"}
%   \def\labSection{e.g. "02"}
% Then % =====================================================================
% Shared preamble for lab sheets.
%
% Caller must \documentclass{article} first, then define:
%   \def\courseTitle{...}
%   \def\labNumber{NN}
%   \def\labTitle{...}
%   \def\labTime{e.g. "30--45 min"}
%   \def\labSection{e.g. "02"}
% Then % =====================================================================
% Shared preamble for lab sheets.
%
% Caller must \documentclass{article} first, then define:
%   \def\courseTitle{...}
%   \def\labNumber{NN}
%   \def\labTitle{...}
%   \def\labTime{e.g. "30--45 min"}
%   \def\labSection{e.g. "02"}
% Then \input{../../template/lab-preamble.tex}
% =====================================================================

\usepackage[utf8]{inputenc}
\usepackage[T1]{fontenc}
\usepackage[english]{babel}
\usepackage[a4paper, margin=2cm]{geometry}
\usepackage{graphicx}
\usepackage{listings}
\usepackage{xcolor}
\usepackage{colortbl}
\usepackage{tikz}
\usepackage{amsmath}
\usepackage{amssymb}
\usepackage{booktabs}
\usepackage{enumitem}
\usepackage{titlesec}
\usepackage{fancyhdr}
\usepackage{hyperref}
\usepackage{tcolorbox}
\tcbuselibrary{skins, breakable}
\usepackage{fontawesome5}

\input{../../../template/tikz-styles.tex}

\providecommand{\courseAuthor}{Industrial Security Lecture Series}
\providecommand{\courseInstitute}{Department of Computer Science}
\providecommand{\companyLogo}{logo}

\definecolor{slideAccent}{HTML}{00205B}
\definecolor{slideSecondary}{HTML}{00A0DC}

% --- Corporate branding overrides ------------------------------------
\IfFileExists{../../../branding.tex}{\input{../../../branding.tex}}{}

\colorlet{labAccent}{slideAccent}
\colorlet{labSec}{slideSecondary}
\definecolor{labCheck}{HTML}{2E7D32}
\definecolor{labWarn}{HTML}{B23A48}

\lstset{
  backgroundcolor=\color{black!4},
  basicstyle=\ttfamily\footnotesize,
  breaklines=true,
  frame=leftline,
  framerule=2pt,
  rulecolor=\color{labAccent},
  numbers=left,
  numberstyle=\tiny\color{black!50},
  showstringspaces=false,
  tabsize=2
}

\setlength{\tabcolsep}{4pt}

% Let TeX add a little extra inter-word stretch as a last resort before
% it reports an Overfull \hbox. Only kicks in on lines that would
% otherwise overflow (long \texttt paths, URLs, CIDRs); never loosens a
% line that already fits. Keeps prose/itemize within the margin.
\setlength{\emergencystretch}{3em}

\pagestyle{fancy}
\fancyhf{}
\renewcommand{\headrulewidth}{0.4pt}
\lhead{\small \courseTitle{} -- Lab \labNumber}
\rhead{\small Maps to Section \labSection}
\cfoot{\thepage}

\newtcolorbox{stepbox}[1]{
  enhanced, breakable, arc=2pt, boxrule=0pt,
  colback=labAccent!7, colframe=labAccent, leftrule=3pt,
  fonttitle=\bfseries\small, coltitle=white,
  colbacktitle=labAccent,
  title={\faTools\ Step #1}
}

\newtcolorbox{warnbox}{
  enhanced, breakable, arc=2pt, boxrule=0pt,
  colback=labWarn!7, colframe=labWarn, leftrule=4pt,
  fonttitle=\bfseries\small, coltitle=white, colbacktitle=labWarn,
  title={\faExclamationTriangle\ Caution}
}

\newtcolorbox{notebox}{
  enhanced, breakable, arc=2pt, boxrule=0pt,
  colback=labAccent!4, colframe=labAccent, leftrule=2pt,
  fonttitle=\bfseries\small, coltitle=white, colbacktitle=labAccent,
  title={\faInfoCircle\ Note}
}

\newtcolorbox{goalbox}{
  enhanced, breakable, arc=2pt, boxrule=0pt,
  colback=labCheck!7, colframe=labCheck, leftrule=3pt,
  fonttitle=\bfseries\small, coltitle=white, colbacktitle=labCheck,
  title={\faBullseye\ Goal}
}

\newtcolorbox{deliverbox}{
  enhanced, breakable, arc=2pt, boxrule=0pt,
  colback=labSec!10, colframe=labSec, leftrule=3pt,
  fonttitle=\bfseries\small, coltitle=white, colbacktitle=labSec,
  title={\faClipboardList\ Deliverable}
}

% --- Solution-sheet boxes (used only by lab-solutions.tex) ----------
\newtcolorbox{solutionbox}[1]{
  enhanced, breakable, arc=2pt, boxrule=0pt,
  colback=labCheck!7, colframe=labCheck, leftrule=3pt,
  fonttitle=\bfseries\small, coltitle=white, colbacktitle=labCheck,
  title={\faCheckCircle\ #1}
}

\newtcolorbox{rubricbox}{
  enhanced, breakable, arc=2pt, boxrule=0pt,
  colback=labSec!7, colframe=labSec, leftrule=3pt,
  fonttitle=\bfseries\small, coltitle=white, colbacktitle=labSec,
  title={\faBalanceScale\ Grading rubric}
}

\titleformat{\section}{\large\bfseries\color{labAccent}}{\thesection}{0.5em}{}

\newcommand{\labheader}{%
  \noindent
  \begin{tikzpicture}
    \fill[labAccent] (0,0) rectangle (\textwidth, -1.6cm);
    \node[anchor=north west, text=white, font=\bfseries\large] at (0.6cm, -0.4cm) {Lab \labNumber{} -- \labTitle};
    \node[anchor=south west, text=white, font=\small] at (0.6cm, -1.3cm) {\courseTitle{} \quad $\bullet$ \quad Maps to Section \labSection{} \quad $\bullet$ \quad \labTime};
    \node[anchor=north east, fill=labSec, text=white, font=\bfseries, inner sep=4pt, rounded corners=2pt]
      at ([xshift=-0.4cm, yshift=-0.4cm]\textwidth,0) {LAB \labNumber};
  \end{tikzpicture}
  \vspace{4mm}
}

% =====================================================================

\usepackage[utf8]{inputenc}
\usepackage[T1]{fontenc}
\usepackage[english]{babel}
\usepackage[a4paper, margin=2cm]{geometry}
\usepackage{graphicx}
\usepackage{listings}
\usepackage{xcolor}
\usepackage{colortbl}
\usepackage{tikz}
\usepackage{amsmath}
\usepackage{amssymb}
\usepackage{booktabs}
\usepackage{enumitem}
\usepackage{titlesec}
\usepackage{fancyhdr}
\usepackage{hyperref}
\usepackage{tcolorbox}
\tcbuselibrary{skins, breakable}
\usepackage{fontawesome5}

% =====================================================================
% Shared TikZ styles for the Industrial Security lecture series.
% Loaded by every slide deck and exercise sheet that uses TikZ.
% =====================================================================

\usetikzlibrary{
  arrows.meta,
  shapes,
  shapes.geometric,
  shapes.symbols,
  shapes.multipart,
  positioning,
  fit,
  calc,
  backgrounds,
  decorations.pathmorphing,
  decorations.pathreplacing,
  decorations.text,
  patterns,
  matrix,
  shadows
}

% --- colour palette --------------------------------------------------
\definecolor{isOT}{HTML}{1F4E79}     % OT (deep blue)
\definecolor{isIT}{HTML}{6F2DA8}     % IT (purple)
\definecolor{isSafety}{HTML}{C00000} % safety red
\definecolor{isNet}{HTML}{2E7D32}    % network green
\definecolor{isAttack}{HTML}{B23A48} % attack red
\definecolor{isAccent}{HTML}{E08E0B} % amber
\definecolor{isMute}{HTML}{6B6B6B}   % grey
\definecolor{isLight}{HTML}{EDF2F8}  % light background
\definecolor{isPLC}{HTML}{2C5282}
\definecolor{isHMI}{HTML}{2F855A}
\definecolor{isSCADA}{HTML}{6B46C1}

% --- generic node styles --------------------------------------------
\tikzset{
  isnode/.style={
    draw, rounded corners=2pt, minimum height=8mm, minimum width=18mm,
    font=\small, align=center, thick
  },
  isbox/.style={isnode, fill=isLight},
  plc/.style={isnode, fill=isPLC!15, draw=isPLC, thick},
  hmi/.style={isnode, fill=isHMI!15, draw=isHMI, thick},
  scada/.style={isnode, fill=isSCADA!15, draw=isSCADA, thick},
  sensor/.style={isnode, fill=isNet!10, draw=isNet, minimum height=6mm, minimum width=14mm, font=\scriptsize},
  actuator/.style={isnode, fill=isAccent!15, draw=isAccent, minimum height=6mm, minimum width=14mm, font=\scriptsize},
  firewall/.style={isnode, fill=isAttack!15, draw=isAttack, thick},
  server/.style={isnode, fill=isMute!15, draw=isMute!80, thick},
  switch/.style={isnode, fill=isLight, draw=black!60, minimum height=6mm, minimum width=14mm, font=\scriptsize},
  workstation/.style={isnode, fill=white, draw=isOT, thick},
  attacker/.style={isnode, fill=isAttack!20, draw=isAttack, thick, font=\small\bfseries},
  inetnode/.style={draw, ellipse, minimum width=24mm, minimum height=10mm, font=\scriptsize, fill=isLight, thick},
  process/.style={isnode, fill=isOT!15, draw=isOT, thick, minimum width=24mm},
  zone/.style={
    draw, rounded corners=4pt, thick, dashed, inner sep=4mm
  },
  conduit/.style={-{Stealth[length=2.5mm]}, thick, isNet!80!black},
  attack/.style={-{Stealth[length=2.5mm]}, thick, isAttack, dashed},
  flow/.style={-{Stealth[length=2.5mm]}, thick, isOT!70!black},
  netline/.style={thick, black!60},
  axisline/.style={-{Stealth[length=2mm]}, thick, black!70},
  tag/.style={font=\scriptsize\itshape, fill=white, inner sep=1pt},
  level/.style={
    draw, fill=isLight, minimum width=0.85\linewidth, minimum height=8mm,
    font=\small, anchor=west
  },
  brace/.style={decorate, decoration={brace, amplitude=4pt, raise=2pt}},
  legendbox/.style={draw, rounded corners=2pt, fill=isLight, inner sep=2mm, font=\scriptsize, align=left}
}

% --- handy macros ----------------------------------------------------
% \plcicon, \hmiicon etc as simple labelled boxes; useful inline.
\newcommand{\plcicon}[1][PLC]{\tikz[baseline=-0.6ex]{\node[plc, minimum width=10mm, minimum height=5mm, font=\scriptsize, inner sep=1pt]{#1};}}
\newcommand{\hmiicon}[1][HMI]{\tikz[baseline=-0.6ex]{\node[hmi, minimum width=10mm, minimum height=5mm, font=\scriptsize, inner sep=1pt]{#1};}}
\newcommand{\fwicon}[1][FW]{\tikz[baseline=-0.6ex]{\node[firewall, minimum width=10mm, minimum height=5mm, font=\scriptsize, inner sep=1pt]{#1};}}


\providecommand{\courseAuthor}{Industrial Security Lecture Series}
\providecommand{\courseInstitute}{Department of Computer Science}
\providecommand{\companyLogo}{logo}

\definecolor{slideAccent}{HTML}{00205B}
\definecolor{slideSecondary}{HTML}{00A0DC}

% --- Corporate branding overrides ------------------------------------
\IfFileExists{../../../branding.tex}{% =====================================================================
% Corporate / institutional branding -- single file to customise the
% look of the entire course (slides, notes, exercises, labs).
%
% HOW TO REBRAND IN UNDER A MINUTE:
%   1. Replace `branding/logo.pdf` with your own logo
%      (PDF preferred; PNG and JPG also work -- name it `logo.<ext>`).
%   2. (Optional) change the two colours below to your brand palette.
%   3. (Optional) override the author/institute lines below.
%   4. Re-run `make` at the repo root. That's it.
%
% Everything in this file has sensible defaults; you can delete a line
% to fall back to the built-in defaults, or comment it out.
%
% This file is loaded automatically by the slides, exercise, and lab
% preambles -- you never need to \input it manually.
% =====================================================================

% --- Logo --------------------------------------------------------------
% Path is resolved from each leaf-build directory; the default points at
% the shared `branding/` folder at the repo root. If you put your logo
% somewhere else, set the full or relative path here (without extension):
\renewcommand{\companyLogo}{../../../branding/logo}

% --- Author / institute (appears on the title page footer) ------------
\renewcommand{\courseAuthor}{}
\renewcommand{\courseInstitute}{}

% --- Brand colours ----------------------------------------------------
% slideAccent      = primary brand colour (title bands, accent rule)
% slideSecondary   = secondary brand colour (section badge, highlight)
% Both use HTML hex codes. Keep enough contrast against white text.
\definecolor{slideAccent}{HTML}{00205B}      % deep navy
\definecolor{slideSecondary}{HTML}{00A0DC}   % light blue accent
}{}

\colorlet{labAccent}{slideAccent}
\colorlet{labSec}{slideSecondary}
\definecolor{labCheck}{HTML}{2E7D32}
\definecolor{labWarn}{HTML}{B23A48}

\lstset{
  backgroundcolor=\color{black!4},
  basicstyle=\ttfamily\footnotesize,
  breaklines=true,
  frame=leftline,
  framerule=2pt,
  rulecolor=\color{labAccent},
  numbers=left,
  numberstyle=\tiny\color{black!50},
  showstringspaces=false,
  tabsize=2
}

\setlength{\tabcolsep}{4pt}

% Let TeX add a little extra inter-word stretch as a last resort before
% it reports an Overfull \hbox. Only kicks in on lines that would
% otherwise overflow (long \texttt paths, URLs, CIDRs); never loosens a
% line that already fits. Keeps prose/itemize within the margin.
\setlength{\emergencystretch}{3em}

\pagestyle{fancy}
\fancyhf{}
\renewcommand{\headrulewidth}{0.4pt}
\lhead{\small \courseTitle{} -- Lab \labNumber}
\rhead{\small Maps to Section \labSection}
\cfoot{\thepage}

\newtcolorbox{stepbox}[1]{
  enhanced, breakable, arc=2pt, boxrule=0pt,
  colback=labAccent!7, colframe=labAccent, leftrule=3pt,
  fonttitle=\bfseries\small, coltitle=white,
  colbacktitle=labAccent,
  title={\faTools\ Step #1}
}

\newtcolorbox{warnbox}{
  enhanced, breakable, arc=2pt, boxrule=0pt,
  colback=labWarn!7, colframe=labWarn, leftrule=4pt,
  fonttitle=\bfseries\small, coltitle=white, colbacktitle=labWarn,
  title={\faExclamationTriangle\ Caution}
}

\newtcolorbox{notebox}{
  enhanced, breakable, arc=2pt, boxrule=0pt,
  colback=labAccent!4, colframe=labAccent, leftrule=2pt,
  fonttitle=\bfseries\small, coltitle=white, colbacktitle=labAccent,
  title={\faInfoCircle\ Note}
}

\newtcolorbox{goalbox}{
  enhanced, breakable, arc=2pt, boxrule=0pt,
  colback=labCheck!7, colframe=labCheck, leftrule=3pt,
  fonttitle=\bfseries\small, coltitle=white, colbacktitle=labCheck,
  title={\faBullseye\ Goal}
}

\newtcolorbox{deliverbox}{
  enhanced, breakable, arc=2pt, boxrule=0pt,
  colback=labSec!10, colframe=labSec, leftrule=3pt,
  fonttitle=\bfseries\small, coltitle=white, colbacktitle=labSec,
  title={\faClipboardList\ Deliverable}
}

% --- Solution-sheet boxes (used only by lab-solutions.tex) ----------
\newtcolorbox{solutionbox}[1]{
  enhanced, breakable, arc=2pt, boxrule=0pt,
  colback=labCheck!7, colframe=labCheck, leftrule=3pt,
  fonttitle=\bfseries\small, coltitle=white, colbacktitle=labCheck,
  title={\faCheckCircle\ #1}
}

\newtcolorbox{rubricbox}{
  enhanced, breakable, arc=2pt, boxrule=0pt,
  colback=labSec!7, colframe=labSec, leftrule=3pt,
  fonttitle=\bfseries\small, coltitle=white, colbacktitle=labSec,
  title={\faBalanceScale\ Grading rubric}
}

\titleformat{\section}{\large\bfseries\color{labAccent}}{\thesection}{0.5em}{}

\newcommand{\labheader}{%
  \noindent
  \begin{tikzpicture}
    \fill[labAccent] (0,0) rectangle (\textwidth, -1.6cm);
    \node[anchor=north west, text=white, font=\bfseries\large] at (0.6cm, -0.4cm) {Lab \labNumber{} -- \labTitle};
    \node[anchor=south west, text=white, font=\small] at (0.6cm, -1.3cm) {\courseTitle{} \quad $\bullet$ \quad Maps to Section \labSection{} \quad $\bullet$ \quad \labTime};
    \node[anchor=north east, fill=labSec, text=white, font=\bfseries, inner sep=4pt, rounded corners=2pt]
      at ([xshift=-0.4cm, yshift=-0.4cm]\textwidth,0) {LAB \labNumber};
  \end{tikzpicture}
  \vspace{4mm}
}

% =====================================================================

\usepackage[utf8]{inputenc}
\usepackage[T1]{fontenc}
\usepackage[english]{babel}
\usepackage[a4paper, margin=2cm]{geometry}
\usepackage{graphicx}
\usepackage{listings}
\usepackage{xcolor}
\usepackage{colortbl}
\usepackage{tikz}
\usepackage{amsmath}
\usepackage{amssymb}
\usepackage{booktabs}
\usepackage{enumitem}
\usepackage{titlesec}
\usepackage{fancyhdr}
\usepackage{hyperref}
\usepackage{tcolorbox}
\tcbuselibrary{skins, breakable}
\usepackage{fontawesome5}

% =====================================================================
% Shared TikZ styles for the Industrial Security lecture series.
% Loaded by every slide deck and exercise sheet that uses TikZ.
% =====================================================================

\usetikzlibrary{
  arrows.meta,
  shapes,
  shapes.geometric,
  shapes.symbols,
  shapes.multipart,
  positioning,
  fit,
  calc,
  backgrounds,
  decorations.pathmorphing,
  decorations.pathreplacing,
  decorations.text,
  patterns,
  matrix,
  shadows
}

% --- colour palette --------------------------------------------------
\definecolor{isOT}{HTML}{1F4E79}     % OT (deep blue)
\definecolor{isIT}{HTML}{6F2DA8}     % IT (purple)
\definecolor{isSafety}{HTML}{C00000} % safety red
\definecolor{isNet}{HTML}{2E7D32}    % network green
\definecolor{isAttack}{HTML}{B23A48} % attack red
\definecolor{isAccent}{HTML}{E08E0B} % amber
\definecolor{isMute}{HTML}{6B6B6B}   % grey
\definecolor{isLight}{HTML}{EDF2F8}  % light background
\definecolor{isPLC}{HTML}{2C5282}
\definecolor{isHMI}{HTML}{2F855A}
\definecolor{isSCADA}{HTML}{6B46C1}

% --- generic node styles --------------------------------------------
\tikzset{
  isnode/.style={
    draw, rounded corners=2pt, minimum height=8mm, minimum width=18mm,
    font=\small, align=center, thick
  },
  isbox/.style={isnode, fill=isLight},
  plc/.style={isnode, fill=isPLC!15, draw=isPLC, thick},
  hmi/.style={isnode, fill=isHMI!15, draw=isHMI, thick},
  scada/.style={isnode, fill=isSCADA!15, draw=isSCADA, thick},
  sensor/.style={isnode, fill=isNet!10, draw=isNet, minimum height=6mm, minimum width=14mm, font=\scriptsize},
  actuator/.style={isnode, fill=isAccent!15, draw=isAccent, minimum height=6mm, minimum width=14mm, font=\scriptsize},
  firewall/.style={isnode, fill=isAttack!15, draw=isAttack, thick},
  server/.style={isnode, fill=isMute!15, draw=isMute!80, thick},
  switch/.style={isnode, fill=isLight, draw=black!60, minimum height=6mm, minimum width=14mm, font=\scriptsize},
  workstation/.style={isnode, fill=white, draw=isOT, thick},
  attacker/.style={isnode, fill=isAttack!20, draw=isAttack, thick, font=\small\bfseries},
  inetnode/.style={draw, ellipse, minimum width=24mm, minimum height=10mm, font=\scriptsize, fill=isLight, thick},
  process/.style={isnode, fill=isOT!15, draw=isOT, thick, minimum width=24mm},
  zone/.style={
    draw, rounded corners=4pt, thick, dashed, inner sep=4mm
  },
  conduit/.style={-{Stealth[length=2.5mm]}, thick, isNet!80!black},
  attack/.style={-{Stealth[length=2.5mm]}, thick, isAttack, dashed},
  flow/.style={-{Stealth[length=2.5mm]}, thick, isOT!70!black},
  netline/.style={thick, black!60},
  axisline/.style={-{Stealth[length=2mm]}, thick, black!70},
  tag/.style={font=\scriptsize\itshape, fill=white, inner sep=1pt},
  level/.style={
    draw, fill=isLight, minimum width=0.85\linewidth, minimum height=8mm,
    font=\small, anchor=west
  },
  brace/.style={decorate, decoration={brace, amplitude=4pt, raise=2pt}},
  legendbox/.style={draw, rounded corners=2pt, fill=isLight, inner sep=2mm, font=\scriptsize, align=left}
}

% --- handy macros ----------------------------------------------------
% \plcicon, \hmiicon etc as simple labelled boxes; useful inline.
\newcommand{\plcicon}[1][PLC]{\tikz[baseline=-0.6ex]{\node[plc, minimum width=10mm, minimum height=5mm, font=\scriptsize, inner sep=1pt]{#1};}}
\newcommand{\hmiicon}[1][HMI]{\tikz[baseline=-0.6ex]{\node[hmi, minimum width=10mm, minimum height=5mm, font=\scriptsize, inner sep=1pt]{#1};}}
\newcommand{\fwicon}[1][FW]{\tikz[baseline=-0.6ex]{\node[firewall, minimum width=10mm, minimum height=5mm, font=\scriptsize, inner sep=1pt]{#1};}}


\providecommand{\courseAuthor}{Industrial Security Lecture Series}
\providecommand{\courseInstitute}{Department of Computer Science}
\providecommand{\companyLogo}{logo}

\definecolor{slideAccent}{HTML}{00205B}
\definecolor{slideSecondary}{HTML}{00A0DC}

% --- Corporate branding overrides ------------------------------------
\IfFileExists{../../../branding.tex}{% =====================================================================
% Corporate / institutional branding -- single file to customise the
% look of the entire course (slides, notes, exercises, labs).
%
% HOW TO REBRAND IN UNDER A MINUTE:
%   1. Replace `branding/logo.pdf` with your own logo
%      (PDF preferred; PNG and JPG also work -- name it `logo.<ext>`).
%   2. (Optional) change the two colours below to your brand palette.
%   3. (Optional) override the author/institute lines below.
%   4. Re-run `make` at the repo root. That's it.
%
% Everything in this file has sensible defaults; you can delete a line
% to fall back to the built-in defaults, or comment it out.
%
% This file is loaded automatically by the slides, exercise, and lab
% preambles -- you never need to \input it manually.
% =====================================================================

% --- Logo --------------------------------------------------------------
% Path is resolved from each leaf-build directory; the default points at
% the shared `branding/` folder at the repo root. If you put your logo
% somewhere else, set the full or relative path here (without extension):
\renewcommand{\companyLogo}{../../../branding/logo}

% --- Author / institute (appears on the title page footer) ------------
\renewcommand{\courseAuthor}{}
\renewcommand{\courseInstitute}{}

% --- Brand colours ----------------------------------------------------
% slideAccent      = primary brand colour (title bands, accent rule)
% slideSecondary   = secondary brand colour (section badge, highlight)
% Both use HTML hex codes. Keep enough contrast against white text.
\definecolor{slideAccent}{HTML}{00205B}      % deep navy
\definecolor{slideSecondary}{HTML}{00A0DC}   % light blue accent
}{}

\colorlet{labAccent}{slideAccent}
\colorlet{labSec}{slideSecondary}
\definecolor{labCheck}{HTML}{2E7D32}
\definecolor{labWarn}{HTML}{B23A48}

\lstset{
  backgroundcolor=\color{black!4},
  basicstyle=\ttfamily\footnotesize,
  breaklines=true,
  frame=leftline,
  framerule=2pt,
  rulecolor=\color{labAccent},
  numbers=left,
  numberstyle=\tiny\color{black!50},
  showstringspaces=false,
  tabsize=2
}

\setlength{\tabcolsep}{4pt}

% Let TeX add a little extra inter-word stretch as a last resort before
% it reports an Overfull \hbox. Only kicks in on lines that would
% otherwise overflow (long \texttt paths, URLs, CIDRs); never loosens a
% line that already fits. Keeps prose/itemize within the margin.
\setlength{\emergencystretch}{3em}

\pagestyle{fancy}
\fancyhf{}
\renewcommand{\headrulewidth}{0.4pt}
\lhead{\small \courseTitle{} -- Lab \labNumber}
\rhead{\small Maps to Section \labSection}
\cfoot{\thepage}

\newtcolorbox{stepbox}[1]{
  enhanced, breakable, arc=2pt, boxrule=0pt,
  colback=labAccent!7, colframe=labAccent, leftrule=3pt,
  fonttitle=\bfseries\small, coltitle=white,
  colbacktitle=labAccent,
  title={\faTools\ Step #1}
}

\newtcolorbox{warnbox}{
  enhanced, breakable, arc=2pt, boxrule=0pt,
  colback=labWarn!7, colframe=labWarn, leftrule=4pt,
  fonttitle=\bfseries\small, coltitle=white, colbacktitle=labWarn,
  title={\faExclamationTriangle\ Caution}
}

\newtcolorbox{notebox}{
  enhanced, breakable, arc=2pt, boxrule=0pt,
  colback=labAccent!4, colframe=labAccent, leftrule=2pt,
  fonttitle=\bfseries\small, coltitle=white, colbacktitle=labAccent,
  title={\faInfoCircle\ Note}
}

\newtcolorbox{goalbox}{
  enhanced, breakable, arc=2pt, boxrule=0pt,
  colback=labCheck!7, colframe=labCheck, leftrule=3pt,
  fonttitle=\bfseries\small, coltitle=white, colbacktitle=labCheck,
  title={\faBullseye\ Goal}
}

\newtcolorbox{deliverbox}{
  enhanced, breakable, arc=2pt, boxrule=0pt,
  colback=labSec!10, colframe=labSec, leftrule=3pt,
  fonttitle=\bfseries\small, coltitle=white, colbacktitle=labSec,
  title={\faClipboardList\ Deliverable}
}

% --- Solution-sheet boxes (used only by lab-solutions.tex) ----------
\newtcolorbox{solutionbox}[1]{
  enhanced, breakable, arc=2pt, boxrule=0pt,
  colback=labCheck!7, colframe=labCheck, leftrule=3pt,
  fonttitle=\bfseries\small, coltitle=white, colbacktitle=labCheck,
  title={\faCheckCircle\ #1}
}

\newtcolorbox{rubricbox}{
  enhanced, breakable, arc=2pt, boxrule=0pt,
  colback=labSec!7, colframe=labSec, leftrule=3pt,
  fonttitle=\bfseries\small, coltitle=white, colbacktitle=labSec,
  title={\faBalanceScale\ Grading rubric}
}

\titleformat{\section}{\large\bfseries\color{labAccent}}{\thesection}{0.5em}{}

\newcommand{\labheader}{%
  \noindent
  \begin{tikzpicture}
    \fill[labAccent] (0,0) rectangle (\textwidth, -1.6cm);
    \node[anchor=north west, text=white, font=\bfseries\large] at (0.6cm, -0.4cm) {Lab \labNumber{} -- \labTitle};
    \node[anchor=south west, text=white, font=\small] at (0.6cm, -1.3cm) {\courseTitle{} \quad $\bullet$ \quad Maps to Section \labSection{} \quad $\bullet$ \quad \labTime};
    \node[anchor=north east, fill=labSec, text=white, font=\bfseries, inner sep=4pt, rounded corners=2pt]
      at ([xshift=-0.4cm, yshift=-0.4cm]\textwidth,0) {LAB \labNumber};
  \end{tikzpicture}
  \vspace{4mm}
}


\begin{document}
\labheader

\begin{goalbox}
Take a greenfield plant from a blank FEED-phase sheet to a defensible
secure-by-design \emph{architecture dossier}. You will derive an IEC
62443-3-2 zone/conduit model with a target Security Level (SL-T) per
zone, specify the Foundational and System Requirements at that target,
draw the network (Purdue + iDMZ + data-diode placement), write a
component procurement spec against the 62443-4-2 Component Requirements,
design a brokered secure-remote-access pattern, embed the SIS without
coupling it to the BPCS, and define the secure-development and
patch-management lifecycle. This is a \emph{design} lab: you produce
architecture artefacts, not exploits. Almost no CLI -- the deliverable
is a dossier an integrator could put out to tender.
\end{goalbox}

\begin{warnbox}
Use the fictional plant defined below only. Do not paste a real client's
IP ranges, asset list, P\&IDs, or vendor names into your dossier, even
partially -- a greenfield design document is exactly the artefact an
attacker most wants to steal. Keep the work to the supplied scenario.
\end{warnbox}

\section*{Scenario -- the greenfield plant}

You are the lead security architect on the FEED (Front-End Engineering
Design) team for \textbf{NorthCell}, a new 12~GWh/year lithium-iron-phosphate
(LFP) battery-cell gigafactory. The plant is being designed from
nothing: there is no legacy equipment, no installed base, no brownfield
excuse. Whatever you specify now becomes contractually binding in the
procurement packages, so design decisions made this week cost a few
engineering hours; the same decisions made after commissioning cost an
outage and 5--10$\times$ the money (INL CCE case studies).

The process, top to bottom:
\begin{itemize}\setlength{\itemsep}{1pt}
  \item \textbf{Electrode coating + drying} -- a continuous web line with
    NMP solvent recovery. A safety-instrumented function (SIF) trips the
    line on solvent-vapour over-concentration (explosion risk).
  \item \textbf{Calendering, slitting, stacking} -- discrete machine cells,
    each with its own PLC and robot.
  \item \textbf{Electrolyte filling + formation} -- the most hazardous
    area; dry-room and inert-atmosphere control; a dedicated SIS.
  \item \textbf{Aging + end-of-line test} -- battery cyclers feeding a
    quality data lake.
  \item \textbf{Plant-wide} -- a process historian, a MES, an energy
    management system (the site draws 60~MW and trades flexibility),
    and a corporate ERP uplink for order/quality data.
\end{itemize}

Drivers: the plant ships globally and is in scope for NIS2 (EU,
\emph{Manufacturing}); the insurer requires an IEC 62443 conformance
statement; the corporate CISO wants a single brokered remote-access
plane for the 14 OEM vendors who will support equipment for 20 years.

\section*{Pre-flight (optional, 10 min)}

\begin{stepbox}{0 -- Tooling}
This lab needs only a drawing surface and a text editor. Use
\texttt{draw.io}, \texttt{excalidraw}, or paper for diagrams; a
spreadsheet or markdown table for the matrices. If you want a reference
topology to react against, the shared lab stack stands in for a single
machine cell:
\begin{lstlisting}[language=bash]
cd master/labs/_stack 2>/dev/null || cd ../../bachelor/labs/_stack
docker compose ps
\end{lstlisting}
Treat one running \texttt{openplc}+\texttt{opcua}+\texttt{historian}
trio as \emph{one} of NorthCell's many machine cells -- it is the unit
you replicate across the floor, not the whole plant. Everything else
in this lab is drawn and written, not typed.
\end{stepbox}

\section*{Step 1 -- Zone \& conduit model + SL-T (IEC 62443-3-2)}

\begin{stepbox}{1 -- Partition the plant and set a target SL per zone}
Run the 62443-3-2 ZCR workflow at \emph{design} fidelity. Define the
System under Consideration as the whole plant network, then partition it
into zones where each zone groups assets that share a threat profile,
a worst-case consequence, and a change-management cadence. A defensible
starting set for NorthCell -- refine it, do not just copy it:
\begin{itemize}\setlength{\itemsep}{1pt}
  \item \textbf{Z-SIS} -- the safety-instrumented systems (coating
    solvent trip, formation/dry-room SIS). Highest consequence: people.
  \item \textbf{Z-CTRL} -- the BPCS: per-cell PLCs, robots, drives,
    line controllers (Purdue L1--L2).
  \item \textbf{Z-SUP} -- supervisory: SCADA/HMI servers, the energy
    management system, the local historian collector (L2--L3).
  \item \textbf{Z-MES} -- manufacturing operations: MES, quality data
    lake, batch records (L3).
  \item \textbf{Z-iDMZ} -- the industrial DMZ brokering all OT$\leftrightarrow$IT
    traffic (L3.5): historian mirror, remote-access broker, update
    repository, reverse proxies.
  \item \textbf{Z-ENT} -- enterprise: ERP, corporate IT, internet
    egress (L4--L5). Out of the OT trust boundary.
\end{itemize}
For \emph{each} zone assign an SL-T (1--4) and justify it in one
sentence anchored to the \emph{worst-case physical consequence}, not to
how interesting the zone looks on a diagram. Then enumerate every
conduit that crosses a zone boundary, one row per conduit, with: name,
source zone, destination zone, protocol, L4 port, direction, data
sensitivity, and the SL-T the conduit must carry (the conduit inherits
the higher of the two zones it joins). Mark which conduits are
data-diode candidates (one-way egress only -- historian/telemetry
export, regulatory feeds).

\medskip
\noindent\textbf{Worked guidance.} SL-T scales with the SL-vector model
of -3-2: estimate unmitigated risk per zone (likelihood $\times$
consequence), compare against the organisation's tolerable risk, and the
gap sets the target. Z-SIS and the formation/dry-room areas are the only
zones where an integrity attack moves \emph{atoms} into a hazard, so they
earn SL-T~3 (a determined, ICS-literate adversary is a realistic threat
for a globally-shipping plant). Z-CTRL/Z-SUP carry SL-T~2--3 depending
on whether the cell can create a physical hazard. Z-iDMZ is high-exposure
but low-consequence-if-breached \emph{provided} segmentation holds, so its
SL-T follows the highest zone it must protect, typically 3. Z-ENT is out
of scope for the OT SL-T but still gets corporate IT controls.
\end{stepbox}

\section*{Step 2 -- FR/SR specification per zone (IEC 62443-3-3)}

\begin{stepbox}{2 -- Turn each SL-T into concrete system requirements}
The seven Foundational Requirements (FR1 IAC, FR2 UC, FR3 SI, FR4 DC,
FR5 RDF, FR6 TRE, FR7 RA) apply to every zone; the SL ladder scales the
\emph{rigor}. Build the FR-per-zone matrix (zones down, FR1--FR7 across,
the target SL in each cell). Then, for your \emph{two} highest-SL zones,
drop one level deeper: pick the System Requirements (SRs) from -3-3 that
the SL-T pulls in, and write what each means at NorthCell.

\medskip
\noindent\textbf{Worked guidance.} The SL ladder inside an FR is
cumulative -- SL3 includes everything in SL1 and SL2 plus the SL3
enhancements. Examples for Z-SIS / Z-CTRL at SL3:
\begin{itemize}\setlength{\itemsep}{1pt}
  \item \textbf{FR1 / SR 1.1, 1.2 (IAC).} Per-user accounts, no shared
    logins; SR~1.5/1.7 enforce hardware-token or MFA for any engineering
    change to a safety-relevant controller.
  \item \textbf{FR2 / SR 2.1 (UC).} Role-based authorisation: only
    \emph{program} role can download logic; operators only acknowledge.
  \item \textbf{FR3 / SR 3.1, 3.4 (SI).} Communication integrity (signed
    or MAC'd control messages where the protocol allows) and software/
    firmware integrity verification on load (SR~3.4).
  \item \textbf{FR4 (DC).} Field protocols rarely need confidentiality;
    document the intentional downgrade and compensate at FR5.
  \item \textbf{FR5 / SR 5.1, 5.2 (RDF).} Enforced zone boundary
    protection; only the conduits listed in Step~1 exist and each is
    filtered (deny-by-default).
  \item \textbf{FR6 / SR 6.1, 6.2 (TRE).} Continuous monitoring;
    auditable event logging forwarded off the zone to a write-once store.
  \item \textbf{FR7 / SR 7.1, 7.3, 7.4 (RA).} DoS resistance, backup
    of control logic, and verified restore -- the cyclers and historian
    must survive a ransomware event without losing batch traceability.
\end{itemize}
Note where you deliberately set an FR below the zone SL-T (e.g.\ FR4 in
the field): a documented, justified downgrade is good engineering; an
undocumented one is a finding.
\end{stepbox}

\section*{Step 3 -- Network design (Purdue, iDMZ, diodes, TSN)}

\begin{stepbox}{3 -- Draw the reference network}
Produce the plant network diagram. Requirements:
\begin{itemize}\setlength{\itemsep}{1pt}
  \item A clear Purdue stratification (L0 field I/O up to L4/5
    enterprise) with the iDMZ at L3.5 as the \emph{only} path between OT
    and IT. No direct L3$\to$L4 conduit may exist.
  \item The iDMZ hosts brokered services in pairs: an inside-facing and
    an outside-facing instance of each (historian mirror, patch repo,
    remote-access broker, reverse proxy). Nothing traverses the iDMZ in
    a single hop -- every flow terminates and is re-originated.
  \item Place at least one \textbf{data diode} where the requirement is
    truly one-way: SIS/safety telemetry and regulatory reporting feeds
    out to Z-MES/Z-ENT. A diode means there is \emph{no} return path --
    no remote control, no ACKs -- so use it only where that is acceptable.
  \item Decide where \textbf{TSN} (Time-Sensitive Networking, IEEE 802.1
    Qbv/Qcc) is justified: the calendering/slitting motion-control cells
    and the coating web line need deterministic, low-jitter cyclic
    traffic. Note that TSN's centralised configuration (CNC/CUC) is
    itself an asset to protect -- a compromised CNC can reshape the
    schedule and stall a line. Aging/test and historian flows do
    \emph{not} need TSN; do not over-specify it.
  \item Segment each machine cell as a micro-zone behind a cell firewall
    (deny-by-default, allow-list of function codes / OPC~UA node sets)
    so a compromised robot cannot reach a neighbouring cell's PLC.
\end{itemize}

\medskip
\noindent\textbf{Worked guidance.} The single most common greenfield
network mistake is a ``flat L2'' control network with one big switch:
cheap to wire, impossible to contain. Spend the cable. Put management
interfaces of switches/firewalls on a dedicated out-of-band management
VLAN reachable only from a hardened jump host -- never in-band on the
control VLAN.
\end{stepbox}

\section*{Step 4 -- Component SL-C + procurement spec (IEC 62443-4-2)}

\begin{stepbox}{4 -- Specify what you will buy, contractually}
A zone reaches its SL-T only if its components can support it. For your
highest-SL zone pick \emph{one} component class (PLC, HMI, or historian)
and write the procurement security specification the buyer will paste
into the RFQ. Map each clause to a -4-2 Component Requirement (CR) and
state the required component capability security level (SL-C).

\medskip
\noindent\textbf{Worked guidance -- a PLC for Z-CTRL at SL-C~3.}
Sample clauses, each tied to a CR:
\begin{itemize}\setlength{\itemsep}{1pt}
  \item CR~1.1/1.2 -- supports individual user accounts and either
    integrates with the plant identity store or supports certificate-based
    auth; no hard-coded or undocumented accounts (CR~1.5).
  \item CR~2.1 -- enforces role-based authorisation for configuration vs.\
    operation; CR~2.8 produces auditable records of every program change.
  \item CR~3.1/3.4 -- verifies the integrity of its own firmware on boot
    and rejects unsigned firmware; supports signed logic download.
  \item CR~3.10 -- supports a secure-by-default configuration delivered
    with all unused services disabled.
  \item CR~7.1/7.6 -- documented DoS resilience and the ability to be
    backed up and restored to a known-good state.
  \item Lifecycle clauses (cross-ref -4-1): SBOM in CycloneDX or SPDX
    with every release; vulnerability handling per ISO/IEC 30111 with a
    disclosure SLA; a 20-year support and security-patch commitment;
    PSIRT contact named in the contract.
\end{itemize}
State explicitly that a vendor unable to deliver an SBOM \emph{today} is
not eligible for a 20-year capital line. Add a short acceptance hook:
the component is rejected on arrival if any clause fails the on-site ATP.
\end{stepbox}

\section*{Step 5 -- Secure remote-access pattern (PAM / jump host)}

\begin{stepbox}{5 -- Design the one brokered access plane}
Fourteen OEMs need remote support for 20 years. Design a single
identity-brokered remote-access pattern that lives in the iDMZ and is the
\emph{only} way any external party touches OT. Specify:
\begin{itemize}\setlength{\itemsep}{1pt}
  \item The flow: vendor $\to$ internet $\to$ SSO+MFA portal (iDMZ
    outside) $\to$ policy/PAM broker $\to$ time-boxed jump host (iDMZ
    inside) $\to$ target asset. No vendor credential ever exists on an
    OT asset; the broker injects a short-lived credential it owns.
  \item Operator-in-the-loop: a named operator approves each session;
    sessions are time-boxed and auto-terminate.
  \item Full session recording (keystroke + screen) to a write-once
    store; the recording is the audit and forensic artefact.
  \item Least privilege per vendor: each OEM reaches only its own
    equipment's micro-zone, never the whole control network.
  \item Break-glass: a documented, dual-control emergency path for when
    the broker itself is down, with mandatory after-the-fact review.
\end{itemize}

\medskip
\noindent\textbf{Worked guidance.} The anti-pattern this kills is the
``vendor VPN straight to the PLC'' -- a flat tunnel that bypasses every
zone boundary you drew in Steps~1--3. The broker enforces FR1/FR2/FR6
(identity, use control, traceability) at one choke point you can audit.
Note the residual risk: the broker is now a crown-jewel asset; it gets
SL-T~3, its own monitoring, and is itself never reachable from the
control zones inbound.
\end{stepbox}

\section*{Step 6 -- Embed safety (IEC 61511 SIS) without coupling to BPCS}

\begin{stepbox}{6 -- Keep safety independent of control and of attackers}
The formation/dry-room and coating areas have safety-instrumented
functions. Design the SIS so that a compromise of the BPCS -- or of the
remote-access plane -- cannot defeat the safety function. Specify:
\begin{itemize}\setlength{\itemsep}{1pt}
  \item \textbf{Independence (61511 separation).} The SIS uses separate
    logic solvers, separate I/O, and separate power from the BPCS.
    A shared HMI may \emph{view} SIS state but must not be able to
    \emph{write} a safety parameter or bypass.
  \item \textbf{No security coupling.} The SIS is its own zone (Z-SIS,
    SL-T~3). The only conduit out is a one-way, diode-protected telemetry
    feed to the historian. There is \emph{no} inbound conduit from the
    BPCS, the iDMZ, or the remote-access broker into Z-SIS.
  \item \textbf{Bypasses and overrides.} Maintenance overrides
    (forces/bypasses) require local, key-switched, dual-authorisation
    and raise an alarm; they can never be set remotely.
  \item \textbf{Security of the safety lifecycle.} Apply 62443 to the SIS
    components too: signed safety logic, change logging, and a
    cyber-security requirements specification feeding the 61511 safety
    requirements specification (the IEC TR 63069 ``functional safety and
    security'' co-engineering view).
\end{itemize}

\medskip
\noindent\textbf{Worked guidance.} The failure to avoid: an integrated
``one HMI controls everything'' design where the same engineering
workstation can download both BPCS logic and SIS logic. That makes the
SIS only as trustworthy as the most-exposed BPCS node -- which defeats
the entire point of an independent protection layer. Draw the SIS as a
visibly separate stack with a single one-way arrow out.
\end{stepbox}

\section*{Step 7 -- Secure-development \& patch-management lifecycle}

\begin{stepbox}{7 -- Define how the design stays secure for 20 years}
A secure-by-design plant decays without a lifecycle. Define two
linked processes:
\begin{itemize}\setlength{\itemsep}{1pt}
  \item \textbf{Vendor SDL gate (IEC 62443-4-1).} What every supplier
    must evidence before and during the contract: SDL conformance,
    SBOM per release, product threat model reviewed each major release,
    vulnerability handling SLA (ISO/IEC 30111), and the on-arrival ATP
    you hooked in Step~4.
  \item \textbf{Asset-owner patch \& maintenance programme (IEC 62443-2-3
    + 2-1).} How NorthCell ingests, tests, and applies patches: a staged
    pipeline (vendor advisory $\to$ risk triage against the zone model
    $\to$ test on an offline twin of the affected cell $\to$ scheduled
    maintenance window $\to$ documented back-out plan). State the
    decision rule for \emph{not} patching (when the compensating control
    in the zone model already reduces the risk below tolerance, and the
    patch risks the process).
  \item \textbf{Backups \& recovery.} The immutable-backup triad
    (production / off-site / WORM) with a quarterly tested restore;
    control logic, HMI projects, historian, and the firewall/switch
    configs are all in scope.
\end{itemize}

\medskip
\noindent\textbf{Worked guidance.} Patch management in OT is not IT's
``patch Tuesday''. The cost of an unplanned trip can exceed the cyber
risk a patch removes, so every patch decision is a risk decision against
the zone model from Step~1, signed off with operations and safety. Write
the back-out plan \emph{before} you touch the running plant.
\end{stepbox}

\section*{Troubleshooting -- common design anti-patterns}

\begin{notebox}
Self-review your dossier against the recurring greenfield mistakes
before you submit:
\begin{itemize}\setlength{\itemsep}{1pt}
  \item \textbf{Flat L2 control network.} One switch, one broadcast
    domain -- cheap to build, impossible to contain. Segment per cell.
  \item \textbf{iDMZ that forwards in one hop.} An iDMZ that lets a flow
    pass through without terminating and re-originating is just a router
    with extra steps. Pair every brokered service inside/outside.
  \item \textbf{Vendor VPN to the PLC.} A flat tunnel bypasses every zone
    boundary. Force all external access through the Step~5 broker.
  \item \textbf{SIS sharing the BPCS network or HMI write path.} Couples
    safety to the most-exposed node; violates 61511 independence.
  \item \textbf{Data diode used where two-way is actually needed.} A diode
    that someone ``works around'' with a side channel is worse than an
    honest firewall. Use diodes only for genuinely one-way flows.
  \item \textbf{SL-T set by exposure, not consequence.} Hardening the
    public-facing zone to SL4 while leaving the safety-adjacent control
    zone at SL2 spends money in the wrong place.
  \item \textbf{Procurement with no security clauses.} If it is not in the
    RFQ it does not exist. Capabilities you cannot buy you cannot deploy.
  \item \textbf{No lifecycle.} A perfect day-one design with no patch or
    backup programme is insecure by year two.
\end{itemize}
\end{notebox}

\section*{Stretch -- threat-model the design (STRIDE / CCE)}

\begin{stepbox}{8 -- Attack your own architecture (stretch)}
Pick one of two lenses and apply it to your dossier:
\begin{itemize}\setlength{\itemsep}{1pt}
  \item \textbf{STRIDE per conduit.} For the three highest-SL conduits,
    enumerate Spoofing, Tampering, Repudiation, Information disclosure,
    Denial of service, Elevation of privilege. For each applicable
    threat, name the control in your design that addresses it and the FR
    it maps to. Any threat with no mapped control is a gap -- list it.
  \item \textbf{INL CCE (Consequence-driven, Cyber-informed Engineering).}
    Phase~1: name the one or two consequences NorthCell cannot accept
    (a dry-room/formation explosion; loss of all batch traceability).
    Phase~2: map the systems-of-systems path an adversary would use to
    cause that consequence. Phase~3: identify the attacker's required
    access. Phase~4: \emph{engineer the consequence out} -- e.g.\ a
    hardwired, non-programmable safety trip that no cyber path can defeat,
    so that even a fully-compromised BPCS cannot cause the explosion.
\end{itemize}
The CCE mindset is the master-level point: do not just add controls --
engineer the worst consequence so it is \emph{physically} unreachable
by any cyber path.
\end{stepbox}

\begin{deliverbox}
Hand in a single PDF -- the \textbf{architecture dossier}:
\begin{enumerate}\setlength{\itemsep}{1pt}
  \item Zone \& conduit model with SL-T per zone and the conduit table,
    diode candidates marked (Step~1).
  \item FR-per-zone matrix plus the SR detail for the two highest-SL
    zones (Step~2).
  \item The reference network diagram: Purdue, iDMZ pairs, diode
    placement, TSN scope, per-cell micro-zones (Step~3).
  \item The component procurement security spec for one component class,
    each clause mapped to a -4-2 CR and an SL-C (Step~4).
  \item The brokered remote-access design with the data flow and the
    break-glass path (Step~5).
  \item The SIS architecture showing independence from the BPCS and the
    single one-way egress (Step~6).
  \item The SDL gate + patch-management + backup lifecycle (Step~7).
  \item (Stretch) the STRIDE-per-conduit table or the CCE consequence
    analysis (Step~8).
\end{enumerate}
A reviewer should be able to read the dossier in fifteen minutes and
reach the same SL-T assignments and zone boundaries you did.
\end{deliverbox}

\section*{Further reading}

\begin{notebox}
\begin{itemize}\setlength{\itemsep}{1pt}
  \item IEC 62443-3-2:2020 -- security risk assessment for system design
    (the ZCR workflow and SL-T derivation).
  \item IEC 62443-3-3:2013 -- system security requirements and SLs
    (FR1--FR7, the SR catalogue).
  \item IEC 62443-4-1:2018 -- secure product development lifecycle.
  \item IEC 62443-4-2:2019 -- technical security requirements for IACS
    components (the CR catalogue and SL-C).
  \item IEC 62443-2-3:2015 / -2-1 -- patch management and the asset-owner
    security programme.
  \item IEC 61511 -- functional safety / SIS lifecycle; IEC TR 63069 --
    co-engineering of functional safety and security.
  \item NIST SP 800-82r3:2023 -- guide to OT security.
  \item INL -- \emph{Countering Cyber Sabotage: Consequence-Driven,
    Cyber-Informed Engineering (CCE)}, Bochman \& Freeman, 2021.
  \item OWASP CycloneDX / SPDX -- SBOM formats; ISO/IEC 30111 -- vuln
    handling.
\end{itemize}
\end{notebox}

\begin{warnbox}
The SL-T values, risk scores, and clause wording in this lab are
pedagogical. A real FEED-phase 62443-3-2 assessment uses the operator's
own risk matrix, the safety classification (SIL) of every SIF, and
walk-downs with the people who will run the plant. A zone-and-conduit
drawing signed without the process-safety engineer in the room is not a
secure-by-design artefact -- it is a liability.
\end{warnbox}

\end{document}
